\begin{otherlanguage}{portuguese}
O monitoramento da biodiversidade por meio de armadilhas fotográficas tornou-se uma ferramenta essencial para estudos ecológicos. No entanto, a análise manual do grande volume de dados gerados por esses dispositivos ainda representa um desafio significativo. Nos últimos anos, métodos baseados em aprendizado profundo têm sido amplamente utilizados para automatizar tarefas como filtragem de imagens vazias, reconhecimento de espécies e identificação de comportamento animal. Apesar desses avanços, tais abordagens geralmente dependem de grandes volumes de dados anotados e apresentam limitações de generalização em cenários reais. Esta proposta tem como eixo central investigar em que condições modelos multimodais de aprendizado de máquina, especialmente modelos baseados em visão e linguagem, podem apoiar a automação de fluxos de trabalho com armadilhas fotográficas em regimes de baixa supervisão. Para isso, o trabalho avalia o potencial desses modelos em tarefas centrais do monitoramento da biodiversidade e examina estratégias para combinar informações visuais e textuais sem depender de treinamento específico para cada tarefa. A primeira etapa do trabalho investiga a capacidade de modelos multimodais, de realizar três tarefas fundamentais: filtragem de imagens vazias, classificação de espécies e identificação de comportamento. Essas tarefas são analisadas em configurações zero-shot e few-shot, considerando também inferência baseada em sequências de imagens. Os resultados indicam que modelos multimodais podem alcançar desempenho competitivo sem treinamento específico para a tarefa, embora sua eficácia varie conforme a tarefa, a arquitetura do modelo e o uso de contexto temporal. Numa segunda etapa, este trabalho propõe e examina uma estratégia de fusão multimodal para classificação de espécies em regime zero-shot. A abordagem combina similaridade visual baseada em CLIP com similaridade textual derivada de descrições geradas por modelos de linguagem. Os resultados mostram que a integração de sinais visuais e textuais pode melhorar a taxa de classificação correta em diferentes cenários, embora os ganhos dependam das características dos dados e do contexto de aplicação. De forma geral, esta proposta contribui com evidências empíricas e metodológicas sobre o uso de modelos multimodais pré-treinados em fluxos de trabalho com armadilhas fotográficas. Em particular, o trabalho demonstra o potencial desses modelos para reduzir a dependência de dados anotados, avalia seus limites em tarefas distintas de monitoramento da biodiversidade e propõe uma estratégia de fusão visual-textual para apoiar a classificação de espécies em cenários zero-shot.

\end{otherlanguage}

% {\setlength{\parindent}{0cm}
% \textbf{Palavras-chave}: Redes neurais profundas, armadilhas fotográficas,
% aprendizagem de máquina, contagem de animais, monitoramento da vida
% selvagem
% }
